\section{风险不是列表，而是估值变量}

本章把每项风险连接到收入、毛利、现金、倍数或行动标签。风险之所以重要，不是因为任何公司都存在风险，而是因为当前价接近基准价值：小幅的执行偏差就足以消耗有限的安全边际。与之相对，长周期技术叙事只有在可审计地转为收入和利润时才增加估值信用。

\begin{exhibitbox}[风险矩阵]
\small
\begin{tabularx}{\textwidth}{L{2.2cm}Y Y Y}
\toprule
风险 & 传导路径 & 领先指标 & 估值/行动响应 \\
\midrule
智驾毛利率 & 芯片、传感器、价格竞争或研发使收入增长不能转为归母 & 分部毛利、成本/收入、费用率 & 下调EPS与P/E；转为高估值风险 \\
订单转收入 & 生命周期年化订单未按项目节奏进入SOP/确认 & 收入增速、项目量产、客户披露 & 不将订单带入基准；低于门槛则下调收入 \\
营运资本 & 备货、账期、客户交付改变现金转换 & 应收、存货、合同负债、CFO、减值 & 加大折现并降低增长可信度 \\
平台/供应链 & SoC路线、供给、成本与认证改变项目经济性 & 平台替代、成本、项目延迟 & 不假设独供；必要时切换熊市情景 \\
OEM与价格竞争 & 客户集中、车型销量、OEM自研或Tier-1打包报价 & 客户/项目变化、毛利、收入解释 & 降低项目持续性与倍数 \\
Physical AI商业化 & 定点、合作或产品发布未转为真实交付、收入、毛利和现金；新业务可能增加研发、库存与应收 & 具名客户、SOP/验收、收入、项目毛利、应收/库存与CFO & 基础EPS维持零信用；无闭环时不因主题抬升倍数 \\
海外制造 & 工厂爬坡、合规、汇率与本地供应链影响现金和利润 & 量产、地区收入/成本、现金流 & 降低海外假设，等待单厂交付证据 \\
\bottomrule
\end{tabularx}
\sourcenote{公司年报、本院风险框架。}
\end{exhibitbox}

最直接的风险是智能驾驶收入增长与毛利质量的错配。2025年智驾收入同比增长32.63\%，但毛利率16.36\%、同比下降3.55个百分点。若该趋势延续，市场可能仍看到高收入增速，而盈利分母和估值倍数会同时承压。模型不以“高阶智驾”自动获得更高利润率，正是为了把这个风险前置。

\section{催化剂与失效条件}

催化剂必须能改变概率，而不仅是增加新闻流。最强的上行证据是正式中报同时显示收入恢复、智能驾驶毛利率稳定和经营现金流匹配；Physical AI的首个强证据是机器人域控实际量产/客户验收或S6付费车队部署，再向收入、项目毛利和回款闭环。单独的订单、合作、定点、POC或工厂进展只提高研究优先级，不改变基础估值。

\begin{exhibitbox}[升级、维持与降级触发]
\small
\begin{tabularx}{\textwidth}{L{2cm}Y Y}
\toprule
状态 & 触发证据 & 行动 \\
\midrule
升级 & 收入恢复中双位数、智驾毛利率企稳、CFO与利润匹配；或新增产品闭合客户--SOP--收入链 & 重新运行盈利桥，评估回撤介入或提高基础锚 \\
维持 & 已确认分部稳定，但新平台仍缺ASP/收入字段 & 维持主业市场支持型观察和Physical AI催化剂跟踪，不把期权加进EPS \\
降级 & 收入持续低增、毛利继续下行、应收/存货压力、项目延期或供应链扰动 & 使用熊市区间，转为高估值风险/观察名单 \\
\bottomrule
\end{tabularx}
\sourcenote{本院监测框架。}
\end{exhibitbox}

反方观点值得被认真对待：若高阶智驾向更多车型扩散、海外客户顺利放量且2027E EPS的实现概率提高，现有17倍左右的前瞻定价可能并不高。我们的反方检验并不否认这种可能，而是要求量产、价格、毛利和现金的证据共同出现。没有这些数据，牛市情景只应是敏感性，不应伪装成基准。

\section{后续研究日历}

报告的有效性到下一次正式财务披露及价格刷新为止。届时应更新市场价格、股本、TTM/前瞻倍数、2026H1收入/利润/现金、分部毛利、应收/库存/合同负债、客户项目进展以及原始券商目标价与方法。若正式披露形成了产品级收入或ASP证据，应先更新客户链审计和增长模型，再更新目标价，不能倒序。

\begin{keyinsight}[投资纪律]
“高质量公司”与“当前价有充足安全边际”是两个不同命题。对德赛西威，本报告认可前者，尚未证明后者。下一次行动改变取决于可审计的经营交付，而不是市场叙事强化。
\end{keyinsight}
